\chapter{结论与展望}

本章凝练全文的核心结论、主要创新点，并简要指出研究的不足与未来方向。

\section{主要结论}

本节按方法、数据和气候效应三个层次总结全文核心发现。

在方法层面，本研究提出的 MPD\_DTW 框架成功实现了全球不同种植制度下水稻生长季的高精度自动识别，兼具空间可迁移性和时间适应性。该框架通过融合物候先验知识与动态时间规整曲线匹配策略，有效解决了全球水稻种植期不固定、生长季长度差异大导致的遥感识别难题 \cite{Dong2016, Maus2016}。基于 MODIS 地表反射率时间序列（MOD09GA），MPD\_DTW 在全球六大洲主要稻区均表现出稳定的识别性能，无需针对特定区域调整物候阈值，显著提升了水稻遥感制图方法的全球普适性。该框架的日尺度处理能力使其能够捕捉年内多季稻的连续种植信息，为水稻复种制度的动态监测提供了新的技术途径。

在数据层面，本研究构建的 GlobalRice500 数据集首次实现了 2001--2022 年全球 500 m 日尺度水稻种植区的长时序动态制图，揭示了全球水稻面积和空间格局的变化特征。GlobalRice500 涵盖全球主要稻区，时间跨度达 22 年，空间分辨率为 500 m，并附带逐像元的生长季起始期（SOS）和结束期（EOS）信息。数据分析表明，2001--2022 年全球水稻面积整体呈波动上升趋势，到 2022 年净增加超过 2100 万公顷；亚洲长期占全球水稻面积的 85\% 以上，非洲是全球水稻面积增长最快的大洲。GlobalRice500 填补了全球覆盖、日尺度、长时序水稻空间分布数据的空白，为全球农业、生态和气候研究提供了重要的基准数据支撑 \cite{Potapov2022, Zabel2019}。

在气候效应层面，本研究基于 GlobalRice500 识别的真实水稻生长季，首次系统揭示了全球稻田地表降温效应的时空格局。研究结果表明，全球稻田在生长季白天具有广泛的局地降温效应，平均降温幅度约为 0.21--0.27 $^{\circ}\mathrm{C}$，且降温强度在分蘖--抽穗期达到峰值。降温强度受稻田比例、田块规模和空间外溢效应的显著调控：稻田比例每增加 1\%，$\Delta$LST 约增加 0.0014 $^{\circ}\mathrm{C}$；大田块（size 11）的平均降温幅度（约 0.69 $^{\circ}\mathrm{C}$）显著高于小田块（size 1，约 0.18 $^{\circ}\mathrm{C}$）；稻田降温效应不仅局限于稻田内部，还向周边非稻田区域产生一定的空间外溢。夜间 $\Delta$LST 约为 0.01--0.02 $^{\circ}\mathrm{C}$，虽显著大于 0 但量级远小于白天，进一步证实了降温机制主要源于白天太阳辐射驱动下的蒸散发过程 \cite{Xue2021RemoteSensing, Weng2025}。

\section{主要创新点}

本节从以下四个方面凝练本论文的主要创新点。

第一，提出融合物候先验知识和曲线匹配策略的全球水稻日尺度识别框架，解决了种植期不固定、生长季长度差异大的全球水稻遥感制图难题。区别于传统的固定物候规则方法和纯数据驱动的机器学习方法，MPD\_DTW 框架将水稻特有的多阶段物候特征嵌入 DTW 曲线匹配过程，实现了对全球不同气候带、不同种植制度下水稻生长季的自动识别，为全球水稻动态制图提供了一套兼具普适性和精度的方法体系 \cite{Dong2016, Xiao2006}。

第二，构建全球长期日尺度 500 m 水稻种植区数据集 GlobalRice500，填补了全球覆盖、日尺度、长时序水稻空间分布数据的空白。GlobalRice500 在时间跨度（2001--2022 年）、空间覆盖（全球主要稻区）、时间分辨率（日尺度）和辅助信息（SOS/EOS）四个维度上均实现了对已有全球水稻数据产品的显著拓展，为水稻物候变化、复种制度转型、耕地扩张过程和气候效应评估等深层次科学研究提供了前所未有的数据基础。

第三，基于真实水稻生长季信息，首次系统量化了全球稻田地表降温效应的时空格局。已有研究对稻田气候效应的关注主要集中于甲烷排放的增温效应，对地表生物物理降温效应的认识仍显不足。本研究利用 GlobalRice500 识别的真实水稻生长季，在全球尺度上系统评估了稻田 LST 相对于周边非稻田的差异，揭示了稻田在生长季白天具有广泛的局地降温效应，并阐明了降温效应的昼夜不对称性和生长季内动态变化规律。

第四，揭示稻田比例、田块规模和周边外溢效应对降温强度的调控作用，深化了对稻田生物物理气候效应空间异质性的认识。本研究在 10 km 网格尺度上量化了稻田比例与降温强度的线性关系，分析了不同田块规模等级下的降温差异，并证实了稻田降温效应向周边非稻田区域的空间外溢。这些发现从景观生态学角度揭示了稻田空间配置对局地气候调节功能的影响机制，为农业景观规划和可持续水稻种植管理提供了科学依据。

\subsection{关键创新}

前述研究背景与文献综述表明，全球稻田遥感监测及其地表生物物理降温效应研究并不是两个彼此孤立的议题。水稻制图为气候效应评估提供空间和时间边界，地表温度效应研究则反过来要求水稻数据不仅能够给出年度稻田范围，还要能够刻画真实生长季、长期变化和稻田空间结构。现有研究已经在区域水稻制图、统计型水稻数据集、灌溉降温和局地稻田小气候效应等方面积累了重要基础，但仍缺乏一条从全球稻田识别、数据产品构建、时空格局分析到稻田地表温度效应评估的完整证据链。因此，本文的研究问题不只是提高某一类分类精度，也不只是报告某一区域的 LST 差异，而是要回答全球水稻这一特殊水—植被—人为管理复合农田系统如何被稳定识别、如何随时间变化，以及这种变化如何影响地表热环境。

第一类创新是全球稻田动态数据缺口问题。已有全球作物或水稻数据产品在全球覆盖、长期连续、空间分辨率和生长季信息之间往往只能满足部分条件。统计分配型产品能够提供宏观面积信息，却难以反映像元尺度稻田边界和真实生长季过程；区域遥感产品空间分辨率较高，但覆盖范围和时间一致性不足；已有全球水稻产品虽然可用于面积统计，却通常缺少日尺度 SOS、EOS、生长季长度和复种制度等动态信息。对于 2001--2022 年全球水稻时空变化分析而言，缺乏稳定可靠、全球一致、日尺度、500 m 分辨率的水稻动态数据集，将直接限制对全球稻田扩张、收缩、区域转移和生长季差异的认识。对于稻田降温效应评估而言，这一数据缺口还会进一步造成 LST 时间窗口错配，使温度差异难以准确归因于水稻生长和稻田水分管理过程。

第二类创新是全球复杂种植制度下的水稻遥感识别方法问题。全球稻田并不是固定月份内出现的单一作物类型，而是受气候背景、灌溉条件、播栽制度、复种制度、品种熟期和农户管理共同控制的动态农田系统。温带单季稻区、亚热带双季稻区、热带多熟稻区、雨养低地稻区和机械化灌溉稻区在水分信号、植被指数曲线、生长季长度和年内周期数量上均存在显著差异。固定阈值和固定物候窗口方法具有可解释性，但难以适应播栽期高度不确定和多季种植并存的全球稻区；监督机器学习和深度学习方法在样本充足区域可以获得较高精度，但跨区域、跨年份和跨管理制度的泛化能力仍受训练样本一致性和长期更新稳定性制约；单纯时间序列匹配方法能够容忍一定物候相位差异，却仍需解决候选生长季搜索、曲线长度差异和多季稻连续识别问题。因此，本文需要回答的核心方法问题是：如何在全球复杂气候区和多样种植制度下，稳定识别不同播栽期、生长周期和复种制度的稻田，并使方法同时具备物候可解释性、时序弹性和长期可迁移性。

第三类创新是基于真实水稻生长季的稻田地表生物物理降温效应问题。已有稻田温度效应研究多证明局地或区域尺度稻田具有降温现象，但全球尺度上稻田相对于周边非稻田农田究竟具有多强的 LST 降低效应、这种效应在白天和夜间是否一致、在水稻生长季内如何随水层和冠层变化而演化、又如何受稻田比例、田块规模和空间邻近关系调控，仍缺乏系统认识。尤其是在全球水稻生长季高度异质的背景下，若采用固定月份或区域平均作物历进行温度比较，就可能将非生长季、轮作期或错季水稻混入分析，从而削弱物理解释。因此，本文需要在真实 SOS--EOS 约束下构建稻田与邻近非稻田农田的对照框架，量化全球稻田地表降温效应的强度、昼夜差异、生长季过程和空间结构调控机制，并明确这一 LST 降温效应与甲烷排放、水资源消耗和近地气温响应之间的边界。

上述三个问题具有递进关系。数据问题回答“全球稻田在哪里、何时生长以及如何变化”，方法问题回答“如何在全球复杂稻作制度下获得可信的水稻动态信息”，效应问题回答“在真实水稻生长季基础上，稻田如何影响地表温度以及这种影响意味着什么”。围绕这一逻辑，本文形成“全球稻田为何难以识别--如何构建普适制图方法--如何形成可信数据产品--全球水稻格局如何变化--稻田如何影响地表温度--这种影响的机制、意义与局限是什么”的研究主线。

\section{不足与展望}

本节与第八章有所呼应，但更简洁，突出未来研究的核心方向。

在数据方面，GlobalRice500 的 500 m 空间分辨率对细碎稻田的刻画能力有限，混合像元效应在散户经营稻区和雨养稻区尤为突出。未来研究可融合 Sentinel-1/2 等 10--30 m 分辨率的多源遥感数据，构建更高空间分辨率的全球水稻数据集，提升对小田块和复杂种植制度的识别精度 \cite{Zhan2021, Han2021}。

在方法方面，MPD\_DTW 框架可向其他具有明显物候阶段的作物（如冬小麦、玉米、大豆等）推广，为全球大宗作物的动态制图提供方法论参考。同时，针对双季稻和三季稻混作区的识别精度仍有优化空间，未来可通过引入年内多峰值检测算法和微波--光学数据融合策略，进一步提升复杂种植制度区域的识别能力。

在气候效应方面，需将地表降温效应与甲烷增温效应纳入统一辐射强迫或 CO$_2$ 当量框架进行综合评估。稻田气候效应具有双重性，只有在统一量纲下完成综合权衡，才能科学回答\"稻田扩张在气候系统上的净效应\"这一关键问题。此外，LST 降温与气温降温之间的定量转换关系也需要通过地面观测和模型模拟进一步厘清。

在应用方面，将 GlobalRice500 嵌入陆面过程模式或气候模型（如 WRF、CESM），评估稻田面积变化、种植北界迁移和种植制度转型对区域能量平衡和气候反馈的影响，是本研究数据产品的重要应用方向。这一工作将为农业气候适应与减缓策略的制定、全球粮食安全评估和农业可持续发展规划提供重要的科学支撑和决策参考。
